\begin{abstract}
Beginners with manual Dobsonian and small alt-azimuth telescopes often cannot find faint objects, because the mount gives no readout of where the tube points. Commercial answers add encoders or a camera dock, and the free phone apps that replace them rely on the phone's compass and gyroscope, which drift and are disturbed by a metal tube. This paper reviews the free and commercial tools in this space and describes AstroFixxer, an open-source Android application that guides a telescope from a phone strapped to it. AstroFixxer extends the AstroHopper web application with a Stellarium-derived offline catalogue of 93,997 deep-sky objects, 8,913 stars and the IAU constellation boundaries, a star-alignment step, and a plate solver written in pure Kotlin that runs on the phone against 579,984 stars to magnitude 10.5 from Gaia DR3 and Hipparcos. The previous sidereal-time formula left pointing errors of up to 801 arcseconds against the astropy reference; the corrected precession model reduces the worst of 66 reference cases to 28 arcseconds. On synthetic star fields the solver solved all 29 test images, including 1-degree eyepiece fields to within 0.5 arcseconds, and refused all 28 negative cases, among them random-star fields. The application has not yet been validated with a real telescope, and the paper states which parts remain to be tested in the field.
\end{abstract}

\begin{IEEEkeywords}
amateur astronomy, push-to navigation, smartphone sensors, plate solving, star identification, offline mobile application
\end{IEEEkeywords}
